\section{Related Work}
\label{sec:related_work}

\paragraph{ViT token pruning.}
DynamicViT~\cite{rao2021dynamicvit}, EViT~\cite{liang2022evit}, A-ViT~\cite{yin2022avit}, Token Merging~\cite{bolya2023tome}, and CrossGET~\cite{shi2024crossget} remove 50--66\% of tokens with under 0.5\,pp ImageNet accuracy loss, but robotic manipulation inverts the spatial-insensitivity assumption. None has been evaluated in closed-loop VLA control.

\paragraph{Token efficiency in VLMs and VLAs.}
FastV~\cite{chen2024fastv}, PruMerge~\cite{shang2024prumerge}, SparseVLM~\cite{zhang2025sparsevlm}, SparseVILA~\cite{khaki2025sparsevila}, and FitPrune~\cite{ye2025fitprune} reduce vision tokens in VLMs using attention or importance signals. VLA-specific methods have emerged: VLA-IAP~\cite{cheng2026vlaiap}, VLA-Cache~\cite{xu2025vlacache}, EfficientVLA~\cite{yang2025efficientvla}, VLA-Pruner~\cite{liu2025vlapruner}, LightVLA~\cite{jiang2026lightvla}, Compressor-VLA~\cite{gao2025compressorvla}, SpecPrune-VLA~\cite{wang2025specprunevla}, and See What Matters~\cite{feng2026seewhatmatters}. Each validates on a single architecture; we ask how architectural choices determine which compression approaches succeed or fail.

\paragraph{Concurrent work.}
Jabbour et al.~\cite{jabbour2025dontrun} independently observe that pruning breaks VLA policies and propose recovery fine-tuning. Our work is complementary: we conduct a systematic cross-architecture analysis identifying the connector as the key factor, with a zero-cost CV diagnostic and characterization of connector--scoring--task interactions.

\paragraph{Vision connectors.}
LLaVA~\cite{liu2023llava} projects patches independently through an MLP; BLIP-2~\cite{li2023blip2} uses cross-attention queries; pixel shuffle~\cite{shi2016subpixel} folds spatial neighborhoods into channels. These create qualitatively different token dependencies whose pruning consequences have not been characterized.
